\documentclass[11pt]{article}

\usepackage[margin=1in]{geometry}
\usepackage{amsmath}
\usepackage{amssymb}
\usepackage{amsthm}
\usepackage{enumitem}
\usepackage{xurl}
\usepackage[colorlinks=true,linkcolor=blue,citecolor=blue,urlcolor=blue]{hyperref}
\setlength{\emergencystretch}{2em}

\theoremstyle{plain}
\newtheorem{theorem}{Theorem}[section]
\newtheorem{lemma}[theorem]{Lemma}
\newtheorem{corollary}[theorem]{Corollary}
\newtheorem{proposition}[theorem]{Proposition}
\theoremstyle{definition}
\newtheorem{definition}[theorem]{Definition}
\newtheorem{hypothesis}[theorem]{Hypothesis}
\theoremstyle{remark}
\newtheorem{remark}[theorem]{Remark}

\newcommand{\M}{\mathcal M}
\newcommand{\W}{\mathcal W}
\newcommand{\V}{\mathcal V}
\newcommand{\QF}{\mathcal Q}
\newcommand{\RR}{\mathbb R}
\newcommand{\status}[1]{\textnormal{\textsf{\footnotesize[#1]}}}
\newcommand{\MachineVerified}{\status{MACHINE-VERIFIED}}
\newcommand{\HumanAudited}{\status{HUMAN-AUDITED}}
\newcommand{\ComputationalEvidence}{\status{COMPUTATIONAL-EVIDENCE}}
\newcommand{\CitedDependency}{\status{CITED-DEPENDENCY}}
\newcommand{\OpenStatus}{\status{OPEN}}
\newcommand{\FailedStatus}{\status{FAILED}}

\title{A Gram Decomposition for the Entropy Kernel in Liu's First Hypothesis}
\author{J.\ Chen\thanks{Repository manuscript package:
\texttt{math/LIU\_H1/paper/}.  Status labels distinguish human proof,
machine-verified algebra, finite computation, and open scope.}}
\date{August 26, 2026}

\begin{document}
\maketitle

\begin{abstract}
\HumanAudited{} We give a self-contained proof of the operational
three-moment formulation of the positive-semidefiniteness hypothesis in
Liu's conditionally-IID approach to the union-closed sets conjecture.  For
\[
  K_{\mathrm L}(s,t)
  =h\bigl((1-s)(1-t)+s(1-s)t(1-t)\bigr),
  \qquad 0\le s,t\le1,
\]
we prove that
\[
  \iint K_{\mathrm L}(s,t)\,d\mu(s)\,d\mu(t)\le0
\]
for every finite signed regular Borel measure \(\mu\) annihilating
\(1\), \(s\), and \(s(1-s)\).  The reduction removes the projected
rank-one terms and leaves a residual kernel \(R\).  An exact Taylor identity
writes \(-R\) as two positive-semidefinite power-series kernels plus an
integral of reciprocal polynomial kernels.  An explicit Lorentz-signature
factorization makes every reciprocal kernel a convergent Gram series.
\MachineVerified{} The finite algebraic identities in this decomposition are
also checked by exact symbolic assertions in the canonical repository script;
those checks corroborate but do not replace the analytic proof.
\OpenStatus{} This manuscript is a candidate resolution of Liu's first
hypothesis only.  Liu's separate Hypothesis~2 and the stronger
union-closed constant derived conditionally from both hypotheses remain open.
A bounded theorem-level literature search found no earlier proof of this
kernel statement; absence from that search is not proof of universal priority.
No claim of peer review, independent acceptance, or end-to-end formal
verification is made.
\end{abstract}

\section{Introduction}
\label{sec:introduction}

\CitedDependency{} A finite family of sets is \emph{union-closed} if it is
closed under pairwise unions.  Frankl's union-closed sets conjecture asks
whether every finite nontrivial union-closed family has an element belonging
to at least half of its member sets; Bruhn and Schaudt survey the history and
classical formulations \cite{bruhn-schaudt}.

\CitedDependency{} Liu introduced a conditionally-IID entropy coupling and a
finite-dimensional reduction for improving quantitative bounds toward that
conjecture \cite{liu}.  The explicit value reported in Liu's Theorem~13 is
conditional on two distinct inputs: the positive-semidefiniteness condition
in Section~V-A and the asserted structure of a nine-parameter global minimizer
in Section~V-B.  We call these inputs Hypothesis~1 and Hypothesis~2,
respectively, following the repository audit terminology.  Liu supported
Section~V-A with the 2,499-node grid code \texttt{frankl3.m} and finite
coefficient code \texttt{frankl7.m}, and Section~V-B with random-start
optimization.  Copies are archived under \path{../literature/code/}; all are
numerical evidence rather than continuum/global proofs.

\HumanAudited{} The purpose of this paper is deliberately narrower than a
new union-closed frequency theorem.  We prove the operational formulation of
Hypothesis~1, including all endpoints and all finite signed measures.  The
argument has three steps:
\begin{enumerate}[label=(\roman*)]
  \item the codimension-three projection is removed exactly;
  \item the resulting residual kernel is written as a sum of three kernels;
  \item the only non-obvious summand is made positive semidefinite by an
        explicit Lorentz-to-Gram conversion.
\end{enumerate}
The residual statement is stronger than the projected assertion because it
requires no moment constraints.

\HumanAudited{} Liu introduced the scale-one kernel and supported its
conditional sign numerically.  The closest rigorous predecessors are the
unperturbed product-entropy kernel of Alweiss--Huang--Sellke
\cite[Lemma~5]{ahs} and Liu's Lemma~11 for sufficiently small perturbation
scale.  Neither reaches \(f(x)=x(1-x)\) at scale one.  We are not aware of an
earlier proof at that endpoint; the full multi-index search cutoff is
2026-08-25, an arXiv-feed refresh on 2026-08-26 found no later collision, and
universal priority remains open.

The proof uses classical closure of PSD kernels under nonnegative sums,
limits, and Schur products \cite[Theorem~7.5.3]{horn-johnson}.  Its reciprocal
kernel is the standard Hilbert-ball kernel
\((1-\langle u,v\rangle)^{-1}\) \cite{agler-mccarthy}.  The
problem-specific content is the exact entropy Taylor regrouping and the
Lorentz factorization placing the normalized polynomial feature map inside
the unit ball.  Search protocol, sources, and limitations:
\path{../LITERATURE_ORIGINALITY.md}.

\paragraph{Status vocabulary.}
Every load-bearing claim is assigned one of the following labels.
\MachineVerified{} means an exact symbolic or validated finite computation
has a replayable code path.  \HumanAudited{} means the continuum mathematical
argument is written and checked as ordinary proof.  \ComputationalEvidence{}
means a finite-precision or finite-dimensional diagnostic only.
\CitedDependency{} identifies material imported from an external source.
\OpenStatus{} marks an unresolved question or incomplete audit, and
\FailedStatus{} marks a route for which an explicit obstruction is recorded.
A machine-verified algebraic identity is not the same as an end-to-end formal
proof in a proof assistant.

\section{Spaces, kernels, and projectors}
\label{sec:setup}

Let \(X=[0,1]\) with its Borel \(\sigma\)-algebra.  Write \(\M(X)\) for the
real vector space of finite signed regular Borel measures on \(X\), equipped
with total-variation norm, and write \(L^2=L^2(X,ds)\) for the real Hilbert
space with Lebesgue measure.  All unqualified integrals below are over
\([0,1]\); double integrals are over \([0,1]^2\).

Define binary entropy in bits by
\begin{equation}
  h(x)=-x\log_2x-(1-x)\log_2(1-x),\qquad 0<x<1,
  \label{eq:entropy}
\end{equation}
with the continuous endpoint values \(h(0)=h(1)=0\).  Put
\begin{equation}
  a(s)=s(1-s),\qquad
  z(s,t)=(1-s)(1-t)+a(s)a(t),\qquad
  K_{\mathrm L}(s,t)=h(z(s,t)).
  \label{eq:liu-kernel}
\end{equation}
The subscript \({\mathrm L}\) distinguishes Liu's entropy kernel from the
auxiliary kernels introduced later.  The functions \(z\) and
\(K_{\mathrm L}\) are symmetric and continuous on the closed square.
Indeed, for fixed \(t\), the function
\((1-s)(1-t)(1+st)\) is nonincreasing in \(s\), so \(0\le z\le1\).

For a continuous symmetric kernel \(F:X^2\to\RR\), define
\begin{equation}
  \QF_F(\mu)=\iint F(s,t)\,d\mu(s)\,d\mu(t),
  \qquad \mu\in\M(X),
  \label{eq:measure-form}
\end{equation}
and define its bounded self-adjoint integral operator on \(L^2\) by
\begin{equation}
  (T_Fg)(s)=\int F(s,t)g(t)\,dt.
  \label{eq:integral-operator}
\end{equation}
We call \(F\) positive semidefinite (PSD) if every finite Gram matrix
\([F(x_i,x_j)]_{i,j=1}^n\) is positive semidefinite; it is negative
semidefinite (NSD) if \(-F\) is PSD.  Lemma~\ref{lem:psd-closure} below connects
this pointwise definition with \eqref{eq:measure-form}.

Set
\begin{equation}
  f_0(s)=1,\qquad f_1(s)=s,\qquad f_2(s)=a(s),
  \qquad
  \W=\operatorname{span}\{f_0,f_1,f_2\}\subset L^2.
  \label{eq:constraint-space}
\end{equation}
The algebraic identity \(a(s)=s-s^2\), together with
\(s^2=s-a(s)\), proves explicitly that
\begin{equation}
  \W=\operatorname{span}\{1,s,s(1-s)\}
    =\operatorname{span}\{1,s,s^2\}.
  \label{eq:span-equality}
\end{equation}
Thus \(\dim\W=3\).  The corresponding moment-annihilating spaces are
\begin{align}
  \V_{\M}
  &=\left\{\mu\in\M(X):
      \int f_j\,d\mu=0\ \text{for }j=0,1,2\right\},
      \label{eq:measure-subspace}\\
  \V_2
  &=\W^\perp
   =\left\{g\in L^2:
      \int f_j(s)g(s)\,ds=0\ \text{for }j=0,1,2\right\}.
      \label{eq:l2-subspace}
\end{align}

For completeness, let
\[
  \boldsymbol f(s)=(f_0(s),f_1(s),f_2(s))^{\mathsf T},
  \qquad
  \Gamma=\int\boldsymbol f(s)\boldsymbol f(s)^{\mathsf T}\,ds
  =\begin{pmatrix}
      1&1/2&1/6\\
      1/2&1/3&1/12\\
      1/6&1/12&1/30
    \end{pmatrix}.
\]
Linear independence makes \(\Gamma\) invertible.  The \(L^2\)-orthogonal
projector onto \(\W\) and the complementary projector onto \(\V_2\) are
\begin{equation}
  (\Pi_{\W}g)(s)
  =\boldsymbol f(s)^{\mathsf T}\Gamma^{-1}
      \int\boldsymbol f(t)g(t)\,dt,
  \qquad
  P_{\V}=I-\Pi_{\W}.
  \label{eq:continuum-projector}
\end{equation}
No other continuum projector is used in the proof.

The associated discrete projector can also be stated without ambiguity.  For
distinct nodes \(s_1,\ldots,s_n\) for which the \(n\times3\) design matrix
\(C\), \(C_{ij}=f_{j-1}(s_i)\), has full column rank, put
\begin{equation}
  P_n=I_n-C(C^{\mathsf T}C)^{-1}C^{\mathsf T}.
  \label{eq:discrete-projector}
\end{equation}
Then \(P_n\) is the Euclidean orthogonal projector onto
\(\ker C^{\mathsf T}\), a subspace of dimension \(n-3\).

\begin{hypothesis}[Operational form of Liu's Hypothesis~1;
\CitedDependency]
\label{hyp:liu-h1}
For every \(\mu\in\V_{\M}\),
\begin{equation}
  \QF_{K_{\mathrm L}}(\mu)
  =\iint K_{\mathrm L}(s,t)\,d\mu(s)\,d\mu(t)\le0.
  \label{eq:liu-h1}
\end{equation}
Equivalently at the level of the named constraints, \(\mu\) annihilates
\(1,s,s(1-s)\), or, by \eqref{eq:span-equality}, \(1,s,s^2\).
\end{hypothesis}

\begin{remark}[The source's codimension count]
\label{rem:codimension}
\CitedDependency{} Liu's Section~V-A calls the affine probability-measure
slice ``codimension 2'' \cite[Section~V-A]{liu}: normalization already fixes
mass, and equations~(48)--(49) add two moments.  Homogeneous signed
perturbations must also annihilate mass, so the accompanying numerical
projector has the three columns \(1,s,s(1-s)\).  By
\eqref{eq:span-equality} these are the independent moments \(1,s,s^2\).
Thus ``codimension two in the probability simplex'' and ``three homogeneous
constraints on signed measures'' describe the same operational tangent space.
\end{remark}

\begin{theorem}[Unprojected residual theorem; \HumanAudited]
\label{thm:residual}
Define the continuous symmetric kernel \(R\) by
\begin{equation}
  R(s,t)
  =A\bigl[B-(1+B)\ln(1+B)\bigr]
   -\bigl[z+(1-z)\ln(1-z)\bigr],
  \label{eq:residual}
\end{equation}
where
\begin{equation}
  u=1-s,\qquad v=1-t,\qquad A=uv,\qquad B=st,
  \qquad z=A(1+B).
  \label{eq:ABz}
\end{equation}
The convention at \(z=1\) is the continuous one.  Then \(R\) is NSD on the
whole closed interval: for every \(\nu\in\M(X)\),
\begin{equation}
  \QF_R(\nu)=\iint R(s,t)\,d\nu(s)\,d\nu(t)\le0.
  \label{eq:residual-nsd}
\end{equation}
No moment, mass, sign, absolute-continuity, or endpoint-exclusion assumption is
placed on \(\nu\).
\end{theorem}

\begin{corollary}[Candidate resolution of Liu's Hypothesis~1;
\HumanAudited]
\label{cor:liu-h1}
Hypothesis~\ref{hyp:liu-h1} holds.  In particular,
\begin{equation}
  \langle g,T_{K_{\mathrm L}}g\rangle_{L^2}\le0
  \quad(g\in\V_2),
  \qquad
  -P_{\V}T_{K_{\mathrm L}}P_{\V}\succeq0
  \quad\text{on }L^2.
  \label{eq:operator-consequence}
\end{equation}
For every discrete node set covered by \eqref{eq:discrete-projector},
\(-P_n[K_{\mathrm L}(s_i,s_j)]_{i,j}P_n\) is positive semidefinite.
\end{corollary}

Theorem~\ref{thm:residual} and Corollary~\ref{cor:liu-h1} are proved in
Sections~\ref{sec:reduction}--\ref{sec:assembly}.

\section{Exact removal of the projector}
\label{sec:reduction}

Use natural logarithms in this section and define
\begin{equation}
  G(x)=x+(1-x)\ln(1-x),\qquad -1\le x<1,
  \quad G(1)=1.
  \label{eq:G-definition}
\end{equation}
On \([0,1]\),
\begin{equation}
  G(x)=\sum_{k=2}^{\infty}\frac{x^k}{k(k-1)}.
  \label{eq:G-series}
\end{equation}

\begin{lemma}[Projected entropy equals the residual form;
\HumanAudited]
\label{lem:reduction}
For every \(\mu\in\V_{\M}\),
\begin{equation}
  (\ln2)\,\QF_{K_{\mathrm L}}(\mu)=\QF_R(\mu).
  \label{eq:exact-reduction}
\end{equation}
All terms are interpreted by continuous extension at \(s=0,1\) and
\(t=0,1\).
\end{lemma}

\begin{proof}
From \(z=uv(1+st)\) and \eqref{eq:entropy},
\begin{align}
  (\ln2)h(z)
  &=-z\ln u-z\ln v-z\ln(1+B)+z-G(z).
  \label{eq:entropy-expansion}
\end{align}
Here, for example, \(u\ln u\) is set to zero at \(u=0\); every apparently
singular product in \eqref{eq:entropy-expansion} then has a continuous value.
Since \(z=A+AB\), move the rank-one term \(AB=a(s)a(t)\) into the remainder.
The result is the pointwise identity
\begin{equation}
  (\ln2)h(z)=-z\ln u-z\ln v+A+R(s,t).
  \label{eq:pointwise-reduction}
\end{equation}
Indeed, the remainder on the right is exactly \eqref{eq:residual}.

It remains to show that the first three kernels on the right of
\eqref{eq:pointwise-reduction} vanish in the quadratic form of \(\mu\).  Their
rank-one decompositions are
\begin{align*}
  -z\ln u
    &=-(u\ln u)v-(us\ln u)(vt),\\
  -z\ln v
    &=-u(v\ln v)-(us)(vt\ln v),\\
  A&=uv.
\end{align*}
The functions \(u=1-s\) and \(us=s(1-s)\) belong to \(\W\); the analogous
statements hold in the \(t\)-variable.  Each rank-one product therefore has at
least one factor annihilated by \(\mu\).  Fubini's theorem
\cite[Chapter~8]{rudin} applies because all extended factors are continuous
and \(\mu\) has finite total variation.
Integrating \eqref{eq:pointwise-reduction} proves
\eqref{eq:exact-reduction}.
\end{proof}

\section{Positive-semidefinite closure and the Taylor identity}
\label{sec:taylor}

\begin{lemma}[Closure rules for continuous PSD kernels;
\HumanAudited]
\label{lem:psd-closure}
Let \(Y\) be a compact metric space.
\begin{enumerate}[label=(\alph*)]
  \item A rank-one kernel \(q(x)q(y)\) is PSD.
  \item A nonnegative linear combination, a pointwise limit, or a pointwise
        product of finite-valued PSD kernels is PSD whenever the resulting
        kernel is finite.  A finite nonnegative integral of PSD kernels is
        PSD whenever its entries are measurable and integrable.
  \item If \(F\) is continuous and PSD on \(Y^2\), then
        \(\iint F\,d\nu\,d\nu\ge0\) for every \(\nu\in\M(Y)\).
\end{enumerate}
\end{lemma}

\begin{proof}
Part (a) follows from
\(\sum_{i,j}c_ic_jq(x_i)q(x_j)=(\sum_i c_iq(x_i))^2\).
For part (b), nonnegative sums and integrals preserve each finite Gram
quadratic form.  Pointwise limits preserve positive semidefiniteness because
the cone of positive semidefinite finite matrices is closed.  Products are
covered by the Schur product theorem \cite[Theorem~7.5.3]{horn-johnson};
equivalently, Gram vectors may be tensor-multiplied.

For part (c), partition \(Y\) into finitely many Borel sets of mesh tending to
zero and choose one point from each nonempty cell.  Replacing \(\nu\) by the
atomic measure that assigns each chosen point the signed mass of its cell
produces nonnegative quadratic forms.  Uniform continuity of \(F\) bounds the
difference from the original double integral by the oscillation of \(F\) on
the product cells times \(\lvert\nu\rvert(Y)^2\), which tends to zero.  The
limit is therefore nonnegative.
\end{proof}

For \(0\le r\le1\), define the polynomial kernel
\begin{equation}
  D_r(s,t)=(1+rB)\,[1-A(1+rB)].
  \label{eq:D-definition}
\end{equation}
It is positive away from \((s,t)=(0,0)\): the first factor is positive, while
\(A(1+rB)\le A(1+B)=z\le1\), with equality only at the corner.

\begin{lemma}[Exact Taylor decomposition; \HumanAudited{} and
\MachineVerified]
\label{lem:taylor}
For every \((s,t)\in[0,1]^2\setminus\{(0,0)\}\),
\begin{equation}
\begin{split}
  -R(s,t)
  ={}&G(A)-AB\ln(1-A)\\
     &+B^2\int_0^1(1-r)\frac{A}{D_r(s,t)}\,dr.
\end{split}
\label{eq:taylor-decomposition}
\end{equation}
At \((s,t)=(0,0)\), assigning the three terms on the right the values
\(1,0,0\), respectively, gives the same identity.  These assignments are
proved continuous in \eqref{eq:log-endpoint} and
Lemma~\ref{lem:reciprocal}.
\end{lemma}

\begin{proof}
Because
\(G(-B)=-B+(1+B)\ln(1+B)\), equation \eqref{eq:residual} gives
\begin{equation}
  -R=G(A(1+B))+A G(-B).
  \label{eq:phi-definition}
\end{equation}
For fixed \(0\le A<1\), let the right-hand side be \(\phi_A(B)\).  Direct
differentiation yields, for all \(B\) in the relevant interval,
\begin{equation}
  \phi_A(0)=G(A),\qquad
  \phi_A'(0)=-A\ln(1-A),\qquad
  \phi_A''(B)=\frac{A}{(1+B)[1-A(1+B)]}.
  \label{eq:phi-derivatives}
\end{equation}
Twice applying the fundamental theorem of calculus gives
\[
  \phi_A(B)-\phi_A(0)-B\phi_A'(0)
  =\int_0^B(B-x)\phi_A''(x)\,dx
  =B^2\int_0^1(1-r)\phi_A''(rB)\,dr.
\]
This is \eqref{eq:taylor-decomposition}.  The excluded value \(A=1\) occurs
only at \((s,t)=(0,0)\), where \(B=0\).  Directly from
\eqref{eq:residual}, \(R(0,0)=-1\), so the assigned endpoint values also give
the identity there.
\end{proof}

The first two kernels in \eqref{eq:taylor-decomposition} are already PSD.
Indeed, by \eqref{eq:G-series},
\begin{align}
  G(A)&=\sum_{k=2}^{\infty}\frac{u^kv^k}{k(k-1)},
  \label{eq:first-psd-series}\\
  -AB\ln(1-A)
  &=\sum_{k=1}^{\infty}\frac{A^{k+1}B}{k}
    =\sum_{k=1}^{\infty}\frac{(u^{k+1}s)(v^{k+1}t)}{k}.
  \label{eq:second-psd-series}
\end{align}
Every summand is rank one with nonnegative coefficient.  At the sole corner
where the logarithmic expression is indeterminate, if
\(m=\max\{s,t\}\), then
\begin{equation}
  0\le AB[-\ln(1-A)]
  \le m^2[-\ln m]\longrightarrow0,
  \label{eq:log-endpoint}
\end{equation}
because \(1-A=s+t-st\ge m\) and \(B=st\le m^2\).  Hence the pointwise-limit
rule in Lemma~\ref{lem:psd-closure} applies on the closed square.

\section{Lorentz factorization and reciprocal Gram kernels}
\label{sec:lorentz}

The remaining integrand in \eqref{eq:taylor-decomposition} contains the kernel
\(AB^2/D_r\).  The numerator is rank one:
\begin{equation}
  AB^2=(1-s)s^2(1-t)t^2=q(s)q(t),
  \qquad q(s)=(1-s)s^2.
  \label{eq:numerator-rank-one}
\end{equation}
The reciprocal denominator is handled next.

\begin{lemma}[Exact Lorentz factorization; \HumanAudited{} and
\MachineVerified]
\label{lem:lorentz}
For \(0\le r\le1\), define
\begin{align*}
  \zeta_0(s)&=-\frac{2rs^2-(r+1)s+1}{r+1},
  &\zeta_1(s)&=1+2rs^2,\\
  \zeta_2(s)&=\frac{s^2(r+2-rs)}{r+2},
  &\zeta_3(s)&=s^3,
\end{align*}
and
\begin{equation}
  \alpha_r(s)=\frac{\zeta_1(s)}{\sqrt{r+1}},\qquad
  \boldsymbol b_r(s)=
  \left(
    \sqrt{r+1}\,\zeta_0(s),
    \sqrt{r(r+2)}\,\zeta_2(s),
    r\sqrt{\frac{2}{r+2}}\,\zeta_3(s)
  \right).
  \label{eq:alpha-b}
\end{equation}
Then, for all \(s,t\in[0,1]\),
\begin{equation}
  D_r(s,t)=\alpha_r(s)\alpha_r(t)
  -\langle\boldsymbol b_r(s),\boldsymbol b_r(t)\rangle_{\RR^3}.
  \label{eq:lorentz-factorization}
\end{equation}
Moreover,
\begin{equation}
  D_r(s,s)\ge s(2-2s+2s^2-s^3)>0
  \qquad(0<s\le1).
  \label{eq:diagonal-lower-bound}
\end{equation}
\end{lemma}

\begin{proof}
With \(\boldsymbol m(s)=(1,s,s^2,s^3)^{\mathsf T}\), direct coefficient
collection in \eqref{eq:D-definition} gives
\begin{equation}
  D_r(s,t)=\boldsymbol m(s)^{\mathsf T}M(r)\boldsymbol m(t),
  \quad
  M(r)=
  \begin{pmatrix}
    0&1&0&0\\
    1&-r-1&2r&0\\
    0&2r&-r(r+2)&r^2\\
    0&0&r^2&-r^2
  \end{pmatrix}.
  \label{eq:coefficient-matrix}
\end{equation}
Let \(S\) be the permutation matrix that swaps the first two coordinates and
put \(\widetilde{\boldsymbol m}(s)=(s,1,s^2,s^3)^{\mathsf T}\).  Exact
multiplication gives
\begin{equation}
  S^{\mathsf T}M(r)S=L(r)\Delta(r)L(r)^{\mathsf T},
  \label{eq:ldl}
\end{equation}
where
\[
  L(r)=
  \begin{pmatrix}
    1&0&0&0\\
    -\frac1{r+1}&1&0&0\\
    -\frac{2r}{r+1}&2r&1&0\\
    0&0&-\frac r{r+2}&1
  \end{pmatrix},
  \qquad
  \Delta(r)=\operatorname{diag}\left(
    -r-1,\frac1{r+1},-r(r+2),-\frac{2r^2}{r+2}
  \right).
\]
A second exact multiplication gives
\[
  L(r)^{\mathsf T}\widetilde{\boldsymbol m}(s)
  =(\zeta_0(s),\zeta_1(s),\zeta_2(s),\zeta_3(s))^{\mathsf T}.
\]
Substitution into \eqref{eq:ldl} is precisely
\eqref{eq:lorentz-factorization}, including \(r=0\), where two negative
coordinates vanish.

For the diagonal bound, set \(q_r=1+rs^2\).  Then
\[
  D_r(s,s)=q_r[1-(1-s)^2q_r],\qquad 1\le q_r\le1+s^2.
\]
The displayed quadratic is concave as a function of \(q_r\), so its minimum
on this interval occurs at an endpoint.  At the endpoints,
\begin{align*}
  1-(1-s)^2&=s(2-s),\\
  1-(1-s)^2(1+s^2)&=s(2-2s+2s^2-s^3).
\end{align*}
The first endpoint value is at least the right side of
\eqref{eq:diagonal-lower-bound}, because
\((2-s)-(2-2s+2s^2-s^3)=s(1-s)^2\ge0\); the second endpoint has the additional
factor \(1+s^2\ge1\).  Finally, the cubic
\(p(s)=2-2s+2s^2-s^3\) decreases from \(2\) to \(1\) on \([0,1]\): its
derivative has discriminant \(-8\) and is negative at \(s=0\).  This proves
\eqref{eq:diagonal-lower-bound}.
\end{proof}

\begin{lemma}[Reciprocal Gram series, including the boundary;
\HumanAudited]
\label{lem:reciprocal}
For each \(r\in[0,1]\), define
\begin{equation}
  H_r(s,t)=
  \begin{cases}
    \displaystyle\frac{AB^2}{D_r(s,t)},&(s,t)\ne(0,0),\\[6pt]
    0,&(s,t)=(0,0).
  \end{cases}
  \label{eq:H-r}
\end{equation}
Then \(H_r\) is a continuous PSD kernel on \([0,1]^2\).  The continuity at
the corner is uniform in \(r\in[0,1]\).
\end{lemma}

\begin{proof}
The function \(\alpha_r(s)>0\) on \([0,1]\).  For \(s>0\), put
\(\boldsymbol c_r(s)=\boldsymbol b_r(s)/\alpha_r(s)\).  From
Lemma~\ref{lem:lorentz},
\[
  1-\|\boldsymbol c_r(s)\|^2
  =\frac{D_r(s,s)}{\alpha_r(s)^2}>0.
\]
Consequently, for \(s,t>0\), Cauchy--Schwarz gives
\(|\langle\boldsymbol c_r(s),\boldsymbol c_r(t)\rangle|<1\), and the geometric
series yields the absolutely convergent identity
\begin{equation}
  \frac1{D_r(s,t)}
  =\frac1{\alpha_r(s)\alpha_r(t)}
   \sum_{k=0}^{\infty}
   \left\langle
     \boldsymbol c_r(s)^{\otimes k},
     \boldsymbol c_r(t)^{\otimes k}
   \right\rangle,
  \label{eq:reciprocal-series}
\end{equation}
where the \(k=0\) inner product is \(1\).  Each summand is a Gram kernel.
Thus \(1/D_r\) is PSD on \((0,1]^2\), and multiplying its feature vectors by
\(q(s)=(1-s)s^2\) shows that \(H_r\) is PSD there.

If exactly one of \(s,t\) is zero, the denominator is positive and the
numerator is zero, so \eqref{eq:H-r} is continuous with value zero.  At the
remaining corner, set \(E_r=1-A(1+rB)\).  Exact factorization gives
\begin{align}
  E_r-s&=(1-s)t\,[1-rs(1-t)]\ge0,\notag\\
  E_r-t&=(1-t)s\,[1-rt(1-s)]\ge0.
  \label{eq:corner-factors}
\end{align}
Thus \(D_r=(1+rB)E_r\ge\max\{s,t\}=:m\), uniformly for
\(r\in[0,1]\).  Since \(0\le AB^2\le s^2t^2\le m^4\),
\[
  0\le H_r(s,t)\le m^3\longrightarrow0
  \qquad\text{as }(s,t)\to(0,0),
\]
again uniformly in \(r\).  For \(\varepsilon\in(0,1)\), define
\[
H_{r,\varepsilon}(s,t)
=H_r\bigl(\max\{s,\varepsilon\},\max\{t,\varepsilon\}\bigr).
\]
This is the pullback of the PSD kernel \(H_r\) on
\([\varepsilon,1]^2\), hence is PSD on the whole closed square.  The continuity
just proved gives \(H_{r,\varepsilon}\to H_r\) pointwise as
\(\varepsilon\downarrow0\), so Lemma~\ref{lem:psd-closure} proves that the
extended \(H_r\) is PSD.
\end{proof}

\section{Assembly of the proof}
\label{sec:assembly}

\begin{proof}[Proof of Theorem~\ref{thm:residual}]
By Lemma~\ref{lem:taylor},
\begin{equation}
  -R=G(A)-AB\ln(1-A)+\int_0^1(1-r)H_r\,dr.
  \label{eq:assembly}
\end{equation}
The first kernel is PSD by \eqref{eq:first-psd-series} and the second by
\eqref{eq:second-psd-series}, with its closed-square extension justified by
\eqref{eq:log-endpoint}.  Lemma~\ref{lem:reciprocal} proves that every \(H_r\)
is continuous and PSD.  Away from \((s,t)=(0,0)\), the map
\((r,s,t)\mapsto H_r(s,t)\) is rational with positive denominator; the uniform
bound \(0\le H_r(s,t)\le\max\{s,t\}^3\) extends it continuously at the corner.
It is therefore bounded and jointly continuous on \([0,1]^3\).  In particular,
every entry \(r\mapsto H_r(x_i,x_j)\) is measurable and integrable, and
integration in \(r\) preserves continuity in \((s,t)\).  Since \(1-r\ge0\),
the integral in \eqref{eq:assembly} is PSD by
Lemma~\ref{lem:psd-closure}.  Hence \(-R\) is a continuous PSD kernel on
\([0,1]^2\).  Applying part (c) of that
lemma to an arbitrary \(\nu\in\M(X)\) gives
\(-\QF_R(\nu)\ge0\), which is \eqref{eq:residual-nsd}.
\end{proof}

\begin{proof}[Proof of Corollary~\ref{cor:liu-h1}]
For \(\mu\in\V_{\M}\), Lemma~\ref{lem:reduction} and
Theorem~\ref{thm:residual} give
\[
  (\ln2)\QF_{K_{\mathrm L}}(\mu)=\QF_R(\mu)\le0.
\]
Since \(\ln2>0\), this is Hypothesis~\ref{hyp:liu-h1}.

If \(g\in\V_2\), the signed measure \(d\mu(s)=g(s)\,ds\) is finite because
\([0,1]\) has finite measure and \(L^2\subset L^1\).  It belongs to
\(\V_{\M}\), so
\(\langle g,T_{K_{\mathrm L}}g\rangle=\QF_{K_{\mathrm L}}(\mu)\le0\).
For arbitrary \(g\in L^2\), apply this inequality to \(P_{\V}g\) to obtain
the operator statement in \eqref{eq:operator-consequence}.

Finally, for nodes \(s_i\) and a vector \(c\in\ker C^{\mathsf T}\), the atomic
measure \(\sum_i c_i\delta_{s_i}\) lies in \(\V_{\M}\).  Thus
\(c^{\mathsf T}[K_{\mathrm L}(s_i,s_j)]c\le0\).  Since \(P_n\) maps onto
\(\ker C^{\mathsf T}\), this is exactly the asserted positive
semidefiniteness of the negative projected matrix.
\end{proof}

\section{Discussion and limits of the result}
\label{sec:discussion}

\HumanAudited{} The proof establishes a stronger unprojected statement for
\(R\), then transfers it to the projected entropy kernel by an exact identity.
No grid limit, spectral truncation, unproved exchange of infinite sums, or
assumption about the support of the signed measure is used.  Three infinite
Gram series appear: \eqref{eq:first-psd-series},
\eqref{eq:second-psd-series}, and \eqref{eq:reciprocal-series}.  The first
converges uniformly because its nonnegative coefficients sum to one.  The
second and third converge absolutely away from their singular corner and
extend to the closed square by \eqref{eq:log-endpoint} and the uniform
\(H_r\le\max\{s,t\}^3\) estimate, respectively.

\OpenStatus{} Liu's Hypothesis~2 is a separate global optimization assertion
for a nine-parameter objective.  Nothing in Theorem~\ref{thm:residual} proves
that assertion, identifies its global minimizers, or validates the stronger
union-closed constant that Liu derives when both hypotheses are assumed.
Accordingly, both Liu's printed decimal \(0.382709087918741\) and the nearby
equation-defined value discussed in the canonical audit remain conditional as
union-closed frequency claims.  This paper does not claim an unconditional
result at either value, nor does it claim Frankl's conjectured constant
\(1/2\).

\HumanAudited{} The completed bounded literature audit found this proof
apparently new through its 2026-08-25 cutoff, with no citation claiming or
proving Liu's Hypothesis~1.  That is a search result, not a universal priority
theorem.  The result remains a candidate until ordinary external mathematical
review; no statement here implies peer review or journal acceptance.

\HumanAudited{} The continuum proof is ordinary mathematics, not a formal
proof-assistant development.  \MachineVerified{} labels below refer only to
specified exact algebraic assertions or validated finite computations.
Therefore the package does not claim end-to-end formal verification.

\appendix
\section{Computational verification and separation of roles}
\label{app:verification}

The analytic proof in Sections~\ref{sec:reduction}--\ref{sec:assembly} is
load-bearing.  The computations in this appendix have narrower roles.  Paths
are relative to this manuscript directory unless explicitly stated otherwise.

\medskip
\noindent\textbf{Artifact:} \path{../verification/independent_exact_checker.py}.
\MachineVerified{} CAS-free coefficient checks over exact rational arithmetic
cover the reduction, projector span, Lorentz factorization, diagonal bound,
and corner factors.

\medskip
\noindent\textbf{Artifact:} \path{../../uc/liu_R_nsd.py}.
\MachineVerified{} A clean isolated run completed successfully and checked the
canonical SymPy identities, finite Arb compression bounds, and floating
diagnostics.  Its exact stdout is
\path{../verification/logs/liu_R_nsd.clean.log}; the analytic continuum
conclusion remains \HumanAudited{}.

\medskip
\noindent\textbf{Artifact:} \path{../../uc/liu_kernel.py}.
\ComputationalEvidence{} This reproduces the projected grid calculation and
the operational three-column projector; it is not a continuum proof.

\medskip
\noindent\textbf{Artifact:} \path{../../uc/liu_moment_cert.py}.
\MachineVerified{} Its exact finite-moment, finite-node Arb, and interval
Cholesky claims replayed successfully; the log is
\path{../verification/logs/liu_moment_cert.clean.log}.
\ComputationalEvidence{} Its floating spectra remain finite-dimensional
corroboration only.

\medskip
\noindent\textbf{Artifact:} \path{../../uc/liu_tail.py}.
\FailedStatus{} This records obstructions to earlier tail-domination routes.
\ComputationalEvidence{} Its sampled diagnostics are not imported into the
proof above.

\subsection{Exact symbolic checks}
\label{app:symbolic}

\MachineVerified{} Two separate exact code paths audit the finite algebraic
core.  The function \path{exact_factorisation_checks()} in
\path{../../uc/liu_R_nsd.py} uses SymPy exact expressions \cite{sympy}.
The package-local
\path{../verification/independent_exact_checker.py} uses only Python integers
and \path{fractions.Fraction}; it does not import the canonical module or any
computer-algebra, numerical, or interval package.  Between them, the exact
assertions cover:
\begin{enumerate}[label=(\arabic*)]
  \item the polynomial reduction \(z=A(1+B)\), the removal \(z-AB=A\), and
        the full rank of \(\{1,s,s(1-s)\}\);
  \item the three values in \eqref{eq:phi-derivatives};
  \item coefficient extraction of \(D_r\) into \(M(r)\), its determinant, and
        its exact characteristic polynomial;
  \item the \(LDL^{\mathsf T}\) and cleared-denominator versions of the Lorentz
        factorization, including \(r=0\);
  \item the diagonal polynomial and nonnegative-factor certificates used in
        \eqref{eq:diagonal-lower-bound};
  \item both exact corner identities in \eqref{eq:corner-factors}.
\end{enumerate}
Successful replay checks transcription and algebra.  Neither code path verifies
Taylor's theorem, logarithm differentiation, kernel closure, convergence,
Fubini's theorem, signed-measure approximation, or the logical assembly of the
result; those are the \HumanAudited{} arguments in the body.

\subsection{Finite validated and numerical corroboration}
\label{app:numerical}

\MachineVerified{} The clean isolated canonical run completed successfully at
the pinned environment in \path{../REPRODUCIBILITY.md}.  The degree
\(8,12,16\) computations give certified negative one-sided upper bounds for
compressed \(R\) eigenvalues through a finite Loewner majorant.  They do not
prove the continuum theorem.  \ComputationalEvidence{} The same run
reconstructs \eqref{eq:taylor-decomposition} at quadrature nodes and inspects
sampled eigenvalues; these are regression tests only.

From the \path{math/} directory, run the independent exact checker with
\begin{quote}\small
\path{python3 -I -B LIU_H1/verification/independent_exact_checker.py}.
\end{quote}
The standard-library-only checker and every canonical clean command, exact
package version, source/log hash, observed exit status, and runtime are recorded
in \path{../REPRODUCIBILITY.md}.  The canonical entry point is executed in an
isolated \texttt{uv} environment rather than trusting the repository virtual
environment.  Exit status zero and the complete labeled output are required.

\subsection{Trust boundary}
\label{app:trust}

The computational trust boundary includes the invoked Python interpreter,
SymPy for exact expression manipulation, and, for interval subroutines, the
Python-FLINT/Arb implementation \cite{arb}.  Floating diagnostics additionally depend on
NumPy and SciPy.  No numerical result is required for the continuum proof.
The manuscript does not assert that these tools constitute a proof assistant
or an independently certified implementation.

\section{Reproducibility checklist}
\label{app:reproducibility}

A reproducibility review should distinguish the following checks.
\begin{enumerate}[label=\textbf{R\arabic*.}]
  \item \HumanAudited{} Verify the definitions
        \eqref{eq:entropy}--\eqref{eq:ABz}, including the base-two versus
        natural-log factor \(\ln2\).
  \item \HumanAudited{} Verify the span identity
        \eqref{eq:span-equality} before comparing the source's codimension
        terminology with its operational projector.
  \item \HumanAudited{} Check every rank-one factor in
        \eqref{eq:pointwise-reduction}; this is the entire use of the three
        moment constraints.
  \item \MachineVerified{} Replay the exact symbolic assertions and compare
        them with \eqref{eq:phi-derivatives}, \eqref{eq:coefficient-matrix},
        and \eqref{eq:ldl}.
  \item \HumanAudited{} Check positivity of \(D_r(s,s)\), absolute convergence
        in \eqref{eq:reciprocal-series}, and the uniform corner estimate in
        Lemma~\ref{lem:reciprocal}.
  \item \HumanAudited{} Check the atomic approximation in
        Lemma~\ref{lem:psd-closure}; this is what extends finite Gram
        positivity to every finite signed regular Borel measure.
  \item \HumanAudited{} Read the bounded literature-search result separately
        from mathematical correctness.  Its 2026-08-25 cutoff and inaccessible
        sources prohibit an unqualified novelty or priority claim.
\end{enumerate}

The LaTeX manuscript itself can be built from this directory with the standard
sequence
\begin{quote}\small
\path{pdflatex main.tex}\quad
\path{bibtex main}\quad
\path{pdflatex main.tex}\quad
\path{pdflatex main.tex}.
\end{quote}
A successful typesetting build checks cross-reference and bibliography syntax;
it does not verify the theorem.

\section{Data and code availability statement}
\label{app:data-code}

No empirical dataset is used in the theorem.  The source of truth for the
proof computations remains the canonical code under \path{../../uc/}; this
paper does not duplicate or fork those scripts.  Package-specific replay
records belong under \path{../verification/}.  The manuscript sources are
\path{main.tex} and \path{refs.bib}.  Availability and reuse are governed by
the repository's existing terms; this paper does not assert an additional
software or data license.

\HumanAudited{} The literature and originality audit at
\path{../LITERATURE_ORIGINALITY.md} is complete through its 2026-08-25 cutoff.
It records three indexed sources as inaccessible and excludes private or
unindexed work; accordingly this paper makes no unqualified priority claim.
Hypothesis~2 code and experiments elsewhere in the repository are outside this
paper's theorem and data claim.

\section*{Acknowledgment of scope}
This manuscript isolates one kernel theorem from a larger union-closed audit.
The status labels are intended to keep analytic proof, symbolic checking,
numerical corroboration, cited context, failed routes, and open questions
logically separate.

\bibliographystyle{plain}
\bibliography{refs}

\end{document}
